% 张三慧《大学物理学》第4版 第二卷第27章《薛定谔方程》
% 章末原题 27.1--27.11；V2 PDF物理页386--387；配套S PDF365--369。
% 与本章另有思考题27.1--27.7区分；图27.14需精确矢量重绘。
\section*{第27章\quad 薛定谔方程}
\subsection*{章末习题27.1--27.11}
\paragraph{习题 27.1}
细胞的长度为 \(a=10^{-5}\,\mathrm m\)，一个粒子质量
\(10^{-14}\,\mathrm g\)。以一维无限深方势阱近似，
计算量子数 \(n_1=100\) 与 \(n_2=101\) 两能级
以及它们的能量差。
\paragraph{习题 27.2}
一氧分子禁闭在宽度为10cm的一维无限深势阱内。
(1) 求基态能量；(2) 若其能量约等于
\(T=300\,\mathrm K\) 时的平均平动能
\(3k_BT/2\)，估计相应量子数 \(n\)，
及 \(n\) 到 \(n+1\) 的能级间距。
\paragraph{习题 27.3}
粒子位于图27.14中两种一维\emph{无限深斜底}势阱
(a)、(b) 内，能量为 \(E_5\) (量子数 \(n=5\))。
依据边界条件、节点数及局部德布罗意波长，
分别画出定态波函数的示意曲线。
\textit{需要对应原图27.14的几何，不能以普通平底井代替。}
\paragraph{习题 27.4}
无限深一维势阱 \((0,a)\) 中的粒子处于基态，
求粒子出现在从某一端起 \(a/4\) 长度区间内的概率。
\paragraph{习题 27.5}
粒子处在一维无限深方势阱中，其本征波函数
按原教材式(27.28)归一化。
求定态 \(n\) 中位置 \(x\) 与 \(x^2\) 的期望值。
\paragraph{习题 27.6}
证明一维无限深势阱中两个不同定态
\(\Psi_m(x,t),\Psi_n(x,t)\) 正交：
\[
\int_0^a \Psi_m^*(x,t)\Psi_n(x,t)\,dx=0
\qquad(m\ne n).
\]
\paragraph{习题 27.7}
一维宽度 \(a\) 的无限深方势阱内粒子
处于能量本征值 \(E_n\) 的定态。
计算粒子对边界壁的平均作用力。
\paragraph{习题 27.8}
一维无限深方势阱内的波函数在边界为零，
等价于两端固定的驻波，要求
井宽容纳德布罗意半波长整数倍。
由此推导能量本征值 \(E_n\)。
\paragraph{习题 27.9}
粒子禁闭在边长 \(a\) 的三维正立方盒子里，
利用三个正交方向的驻波量子数
\(n_x,n_y,n_z=1,2,\dots\)
推导粒子的能量谱。
\paragraph{习题 27.10}
谐振子基态的定态空间波函数为
\(\psi_0(x)=Ae^{-\alpha x^2}\)，
其中 \(A,\alpha\) 为常量。
代入一维谐振子薛定谔方程并利用该等式对所有
\(x\) 成立，求证零点能 \(E_0=h\nu/2\)。
\paragraph{习题 27.11}
将氢分子的振动近似为简谐振子，
等效劲度 \(k=1.13\times10^3\,\mathrm{N/m}\)、
质量 \(m=1.67\times10^{-27}\,\mathrm{kg}\)。
求其本征能级（单位eV），
及相邻能级辐射光子的能量与波长。
